% ==============================================================================
% Fichier     : main.tex
% Titre       : Résumé académique -- Processus Investigatif Universel
% Institution : Université de Yaoundé I -- ENSP
% ==============================================================================

\documentclass[12pt, a4paper]{report}

\usepackage{style}
\usepackage{multicol}
\sloppy

% ==============================================================================
% VARIABLES
% ==============================================================================
\newcommand{\Universite}{Université de Yaoundé I}
\newcommand{\Etablissement}{École Nationale Supérieure Polytechnique}
\newcommand{\Departement}{Département de Génie Informatique}
\newcommand{\AnneeAcademique}{2025 -- 2026}
\newcommand{\LogoEtab}{logo_polytech.png}
\newcommand{\CodeUE}{GES4062}
\newcommand{\TitreRapport}{TECHNIQUES D'INVESTIGATION NUMERIQUE AVANCEE}
\newcommand{\ThemeProjet}{\LARGE Résumé de la méthodologie du cours d'Investigation Numérique}
\newcommand{\Etudiant}{MAKEU TENKU STELY BELVA}
\newcommand{\Matricule}{24P826}
\newcommand{\Filiere}{Cybersécurité et Investigation Numérique}
\newcommand{\Niveau}{Niveau 4}
\newcommand{\Encadreur}{Ing. THIERRY MINKA}
\newcommand{\GradeEncadreur}{Enseignant -- ENSP}
\newcommand{\DateSoumission}{Mars 2026}

% ==============================================================================
\begin{document}
	
	% PAGE DE GARDE
	\input{PageGarde}
	
	% ==============================================================================
	% PAGE 1
	% ==============================================================================
	\thispagestyle{plain}
	\vspace*{0.2cm}
	
	% En-tête du résumé
	
	
	\vspace{0.35cm}
	
	% ---------------------------------------------------------------
	% INTRODUCTION
	% ---------------------------------------------------------------
	\noindent
	Le \termecle{Processus Investigatif Universel} (PIU) constitue un cadre
	méthodologique structuré, conçu pour conduire toute investigation --
	individuelle, organisationnelle ou judiciaire -- de manière rigoureuse,
	traçable et juridiquement défendable. Son architecture modulaire repose
	sur \textbf{cinq phases séquentielles}, \textbf{dix-sept étapes opérationnelles}
	et \textbf{trois points de contrôle qualité} (Quality Gates), chaque étape
	étant déclinable selon trois niveaux de profondeur : minimal, standard
	et avancé, calibrés en fonction de la criticité du cas.
	
	\vspace{0.4cm}
	
	% ---------------------------------------------------------------
	% SECTION 1
	% ---------------------------------------------------------------
	\begin{boxobjectifs}
		\textbf{Objectifs fondamentaux du PIU :}
		\begin{itemize}
			\item Établissement factuel objectif et reproductible
			\item Préservation de l'intégrité probatoire à chaque étape
			\item Traçabilité complète permettant l'audit indépendant
			\item Production de conclusions résistantes à la contestation contradictoire
		\end{itemize}
	\end{boxobjectifs}
	
	\vspace{0.35cm}
	
	% ---------------------------------------------------------------
	% SECTION 2 : LES 5 PHASES
	% ---------------------------------------------------------------
	{\color{CP}\large\sffamily\bfseries\faAngleDoubleRight\enspace
		Architecture du Processus en Cinq Phases}
	\vspace{0.15cm}
	\color{CV}\rule{\linewidth}{0.8pt}
	\color{CT}
	
	\vspace{0.3cm}
	
	\textbf{Phase 1 -- Initialisation (E1--E3).}
	Cette phase établit la recevabilité de la demande investigative. Elle
	comprend l'évaluation préliminaire de la demande (E1), la décision
	formelle d'engagement matérialisée par un mandat d'investigation (E2),
	et la planification opérationnelle (E3). Le \termecle{Quality Gate 1}
	valide la disponibilité des ressources, la clarté du périmètre et
	l'opérationnalité des outils avant tout démarrage effectif.
	
	\vspace{0.25cm}
	
	\textbf{Phase 2 -- Acquisition et préservation (E4--E6).}
	Cette phase constitue la base probatoire de l'investigation. La collecte
	systématique des éléments (E4) s'effectue selon des protocoles documentés,
	incluant horodatage et identification de l'opérateur. La sécurisation (E5)
	applique le principe 3-2-1 (trois copies, deux supports, une localisation
	distincte) avec calcul d'empreintes cryptographiques (SHA-256 minimum).
	La cartographie des parties prenantes (E6) établit la matrice
	relationnelle complète. Le \termecle{Quality Gate 2} certifie l'intégrité
	et la recevabilité juridique du corpus probatoire.
	
	\vspace{0.25cm}
	
	\textbf{Phase 3 -- Investigation analytique (E7--E10).}
	Phase centrale du processus, elle mobilise la collecte testimoniale (E7),
	la génération d'au moins cinq hypothèses distinctes (E8) -- incluant
	obligatoirement une hypothèse nulle et des hypothèses improbables --,
	la confrontation systématique hypothèses/preuves via une matrice de
	compatibilité (E9), la \termecle{falsification active} inspirée du
	principe épistémologique de Popper (E9bis), et la validation
	multi-méthodes des conclusions (E10). Le \termecle{Quality Gate 3}
	conditionne la progression ; les itérations de cette phase (retours
	vers E4--E6 ou E7--E10) constituent une caractéristique normale et
	souhaitable d'une investigation rigoureuse.
	
	\vspace{0.25cm}
	
	\textbf{Phase 4 -- Formalisation et transmission (E11--E13).}
	Le dossier consolidé s'organise en cinq volumes (synthèse, analyses,
	preuves, témoignages, annexes). Le rapport final adopte une structure
	tripartite : synthèse exécutive (2--4 pages), rapport technique
	(20--100 pages), dossier probatoire. La transmission s'effectue de
	manière sécurisée et traçable, avec bordereau signé et accusé de
	réception contresigné.
	
	\vspace{0.25cm}
	
	\textbf{Phase 5 -- Suivi post-transmission (E14--E17).}
	Cette phase, essentielle en contexte contentieux, couvre la préparation
	à la défense contradictoire, la maintenance probatoire, l'analyse des
	contre-expertises et la clôture avec retour d'expérience (REX)
	documentant les enseignements pour les investigations futures.
	
	
	
\end{document}